\begin{otherlanguage}{english}
\chapter*{\centering Abstract}
Demand heterogeneity makes replenishment-policy selection a decision in its own right, alongside parameter calibration. This dissertation proposes the Adaptive Inventory Policy Engine (AIPE), a modular methodology supporting logistics and inventory managers in reducing avoidable costs while preserving service. The procedure comprises five stages: specifying the decision, representing context, building comparable evidence, recommending a policy, and evaluating its consequences. Its Policy Selection Engine (PSE) maps profiles to policy families or estimates recommendations for each series, with an intervention rule that can retain the reference policy. We evaluate seven policies in two separate applications. In Bahia, 145 series at two origins yield 4,350 fits; profile-based selection increases mean cost by 1.93\% and mean fill rate by 3.89 percentage points relative to global selection. In the M5 pilot, 30 series from ten stores at three origins yield 1,350 fits. The class-based predictive selector has higher cost and lower service. An exploratory extension based on predicted cost and service, with fallback, reduces mean cost by 1.56\% and increases service by 0.27 percentage points, but changes only three decisions; one series accounts for approximately 92\% of savings. Retrospective oracles distinguish selection errors, portfolio gaps, and physical infeasibility. The contribution is a procedure linking each recommendation to evidence to support human decisions. Practical relevance depends on context; evaluated costs are scenario-based. The policy family remains fixed within each seven-cycle block, without validating cycle-by-cycle PSE revision. The results demonstrate execution through the evaluated interfaces and delimit the utility of the implementations. Retrospective cohorts, different calibrations, short horizons, and development on the pilot restrict generalization and require independent confirmation.

\vspace{2em}
\textbf{Keywords}: Selection methodology; Inventory management; Contextual policy selection; Service level; AIPE.
\end{otherlanguage}
